\begin{theorem}[The component along $\Delta$]
\label{thm:c9-D1}
Let $\hbar F_1^{(0)}$ contain the finite counterterm of Theorem \ref{thm:c4-finite}, let $\hbar F_1^{(1)}(\Delta)$ solve the quadratic part of \eqref{eq:c9-arity1}, with the part \eqref{eq:c9-Fperp} in the sector $(A|_{W_{\mathbb R}^\perp},A^*|_{W_{\mathbb R}^\perp})$, and let $\psi=0$. For a scalar bilinear $O$ the arity-one counterterm of $O$ along $\Delta$ vanishes, and on the part with two legs $A|_{W_{\mathbb R}^\perp}$ the local representative along $\Delta$ exists and is given by the master formula of Theorem \ref{thm:c11-classes} (iii),
\[
D_1(O)=\Upsilon_\Delta(A|_{W_{\mathbb R}^\perp},A|_{W_{\mathbb R}^\perp})-\big\{\hbar F_1^{(1)}(\Delta),O\big\},
\]
with the two-body kernel $\Upsilon_\Delta$, the coefficient of $\varepsilon^0$ of the defect cycles of the pair of letters of $O$ along $\Delta$, equal to $\tfrac{44}3\,O-\tfrac13\{S^{(1)}(\Delta),O\}$ units on the traceless part and $-\tfrac{100}3\,O-\tfrac13\{S^{(1)}(\Delta),O\}$ units on the trace. Hence
\begin{equation}
\begin{aligned}
D_1(O)&=-\big(\tfrac13+\theta\big)\frac{\hbar\operatorname{tr}_{\mathfrak g}(\mathrm{ad}_X\mathrm{ad}_Y)}{16\pi^2(X,Y)}\,\{S^{(1)}(\Delta),O\}&&\text{for the traceless part},\\
D_1(O)&=-48\,\frac{\hbar\operatorname{tr}_{\mathfrak g}(\mathrm{ad}_X\mathrm{ad}_Y)}{16\pi^2(X,Y)}\,O-\big(\tfrac13+\theta\big)\frac{\hbar\operatorname{tr}_{\mathfrak g}(\mathrm{ad}_X\mathrm{ad}_Y)}{16\pi^2(X,Y)}\,\{S^{(1)}(\Delta),O\}&&\text{for the trace},
\end{aligned}
\label{eq:c9-D1}
\end{equation}
$\{S^{(1)}(\Delta),O\}=2(E+2)O$: twice the coefficient of Theorem \ref{thm:c6-dil} times $O$, plus $-(\tfrac13+\theta)$ units times the classical transformation. At the value $\theta=-\tfrac13$ fixed by the arity-two equations (Theorem \ref{thm:c9-theta2}),
\[
D_1(O)=0\quad\text{(traceless part)},\qquad D_1(O)=-48\,\text{units}\cdot O=\frac{3g^2N}{2\pi^2}\,O\quad\text{(trace, }\mathfrak g=\mathfrak{su}(N)\text{)},
\]
and these values are the scheme-independent content: they are unchanged by every free direction of the action-side solution compatible with the arity-two equations, by the finite parts $\psi$ of Lemma \ref{lem:c10-psi} and by the finite counterterms (Theorem \ref{thm:c8-invariants} (iii)). The contributions are those of the following table.
\end{theorem}

\begin{tabular}{l|c|c}
contribution, part with two legs $A|_{W_{\mathbb R}^\perp}$ & traceless part & trace \\
\hline
defect cycles of the pair along $\Delta$, $\Upsilon_\Delta$ & $\tfrac{44}3\,O-\tfrac13\{S^{(1)}(\Delta),O\}$ & $-\tfrac{100}3\,O-\tfrac13\{S^{(1)}(\Delta),O\}$ \\
$-\{\hbar F_1^{(1)}(\Delta),O\}$ & $-\tfrac{44}3\,O-\theta\,\{S^{(1)}(\Delta),O\}$ & $-\tfrac{44}3\,O-\theta\,\{S^{(1)}(\Delta),O\}$ \\
$-\{I^{CT}_{\mathfrak{psl}(4|4)}[\varepsilon]^{(1)}(\Delta),O\}$, arity-one counterterm of $O$ & $0$ & $0$ \\
\hline
$D_1(O)$ & $-\big(\tfrac13+\theta\big)\{S^{(1)}(\Delta),O\}$ & $-48\,O-\big(\tfrac13+\theta\big)\{S^{(1)}(\Delta),O\}$ \\
\end{tabular}
\tabcaption{The component $D_1(O)$ of a scalar bilinear on the part with two legs $A|_{W_{\mathbb R}^\perp}$, in units; for the trace, $O$ stands for $\int\omega\,(A|_{W_{\mathbb R}^\perp},A|_{W_{\mathbb R}^\perp})$. At the point of the density the first row is $\tfrac{40}3\,O$ and $-\tfrac{104}3\,O$. By Corollary \ref{cor:c11-DPdefect} it equals $4$ times the coefficient of $\log\varepsilon$ in $\partial_\delta I^{CT}[\varepsilon]$ (Theorem \ref{thm:c5-Oeps}), $16\,O$ and $-32\,O$, plus the composite pair kernel $-\tfrac43O-\tfrac13\{S^{(1)}(\Delta),O\}$, an identity of two evaluations of the same flow.}

\begin{proof}
By Theorem \ref{thm:c11-classes} (iii) at $T=\infty$, which applies along $\Delta$ (Proposition \ref{prop:c8-Tsafe}), $D_1(O)$ is the two-body kernel of the single pair of letters of $O$ minus the bracket of $O$ with $\hbar F_1^{(1)}(\Delta)$. The two-body kernel on the outputs $A|_{W_{\mathbb R}^\perp}A|_{W_{\mathbb R}^\perp}$ is the coefficient of $\varepsilon^0$ of the defect cycles of the pair: cycles through the two letters with one special edge, vertices of $I$ and the vertex $S^{(1)}(\Delta)$ at the special edge, with one leg on it and the other on an edge $P$ or at a letter, or on a line, the composite kernel; they carry no $\log\varepsilon$ (Theorem \ref{thm:c11-oddsimplex} (i)), their poles cancel against the poles of the counterterm trees (\eqref{eq:c11-pole}), and their exact evaluation, with the position realised by a momentum injected at $S^{(1)}(\Delta)$, gives the first row, whose position part along $b$ is $2\,\Upsilon_{P_b}=-\tfrac23$ units times $\partial_bO$ (Theorem \ref{thm:c9-P1} (ii)). The graphs with a special edge on a tree tend to $0$ together with the special-edge parts of their counterterm partners, and the graphs whose cycle passes through vertices of the action alone tend to $0$ by the arity-zero equation, whose two-leg part vanishes in the sector $(c,A|_V)$ (Theorem \ref{thm:c4-finite}), and by \eqref{eq:c9-arity1} on two legs. The one-body term is computed by Lemma \ref{lem:c4-count} and Lemma \ref{lem:c9-SDelta} (i): $-\tfrac{22}3\cdot2\,O-\theta\{S^{(1)}(\Delta),O\}$ units from \eqref{eq:c9-Fperp}, and $0$ from the counterterm, whose quadratic part contains no $A^*|_{W_{\mathbb R}^\perp}$ (Theorem \ref{thm:c9-quadratic} (i)). The arity-one counterterm of $O$ along $\Delta$ is a colour-singlet local functional of degree $-1$ and Weyl weight $2$, with terms linear in the position carrying densities of weight $3$ (the weight counting in the proof of Proposition \ref{prop:c7-select}); every colour singlet of degree $-1$ contains an antifield and a second adjoint field and has weight at least $4$, so it vanishes. The sum is the last row. By Theorem \ref{thm:c5-Oeps} and Lemma \ref{lem:c6-count}, $4$ times the coefficient of $\log\varepsilon$ in $\partial_\delta I^{CT}[\varepsilon]$ minus $16\,O$ is twice the coefficient of Theorem \ref{thm:c6-dil} times $O$: $2\cdot0$ for the traceless part and $2\cdot(-24)$ for the trace. The invariance statement: by Theorem \ref{thm:c9-theta2} the coefficient $\theta$ is the same in every solution of the arity-two equations, the inner and outer directions leave the uniform coefficient along $\Delta$ unchanged, $\lambda_\Delta=B_\Delta=0$, and Lemma \ref{lem:c10-psi} and Theorem \ref{thm:c8-invariants} (iii) give the rest.
\end{proof}
